\chapter[Band 8: Bijektive Funktionen]{Band 8 auf einen Blick}
\noindent\textbf{Bijektive Funktionen}\par\medskip
\noindent\emph{Lesart.} Bijektionen verbinden die beiden getrennt
entwickelten Axiome der Injektivität und Surjektivität. \(F^{-1}\)
bezeichnet hier die Umkehrfunktion einer Bijektion; \(F^{-1}[T]\)
bleibt die Urbildmenge.
\begin{BandUebersicht}
\textbf{Bijektive Funktion}
\UebFormel{\begin{aligned}
&F\colon A\bij B:\quad F\colon A\to B,\\
&F\colon A\inj B,\qquad F\colon A\sur B.
\end{aligned}}
\UebQuellen{\FormulaRefAuto{Bijektive Funktion}[def]}
&
\textit{Komposition von Bijektionen}
\UebFormel{\begin{aligned}
&F\colon A\bij B\dsep G\colon B\bij C\\
&\hspace{12mm}\vdash G\circ F\colon A\bij C.
\end{aligned}}
\UebQuellen{\FormulaRefAuto{BijectiveFunctionComposition}[thm]}\\
\addlinespace[.8ex]
\textbf{Identität und Umkehrfunktion}
\(\Id_A\) ordnet jedem \(x\in A\) sich selbst zu. Für
\(F\colon A\bij B\) wird definiert:
\UebFormel{F^{-1}(y)\coloneqq
\iota x\,(x\in A\land F(x)=y)\quad(y\in B).}
\UebQuellen{\FormulaRefAuto{InverseFunctionDef}[def]}
&
\textit{Beidseitige Umkehrung}
\UebFormel{\begin{aligned}
&F\colon A\bij B\vdash F^{-1}\colon B\bij A,\\
&F\colon A\bij B\dsep x\in A\vdash{}\\[-2pt]
&\qquad F^{-1}(F(x))=x,\\
&F\colon A\bij B\dsep y\in B\vdash{}\\[-2pt]
&\qquad F(F^{-1}(y))=y.
\end{aligned}}
\UebQuellen{\FormulaRefAuto{InverseFunctionBijection}[thm];
\FormulaRefAuto{InverseFunctionLeftInverse}[thm];
\FormulaRefAuto{InverseFunctionRightInverse}[thm]}\\
\addlinespace[.8ex]
\textbf{Bijektive Einschränkungen}
Sei \(F\colon A\bij B\), \(S\subseteq A\), \(T\subseteq B\).
Die Einschränkung verwendet weiterhin die Werte von \(F\).
\UebQuellen{\FormulaRefAuto{RestrictionDef}[def]}
&
\textit{Punktweises Kriterium}
\UebFormel{\begin{aligned}
&F\restriction_S\colon S\bij T\ \leftrightarrow\\
&\quad\forall x\in A\,(x\in S\leftrightarrow F(x)\in T),\\
&F\restriction_{F^{-1}[T]}\colon F^{-1}[T]\bij T.
\end{aligned}}
\UebQuellen{\FormulaRefAuto{BijectiveRestrictionIffPointwiseCompatibility}[thm];
\FormulaRefAuto{F\colon A\bij B \dsep T\subseteq B \vdash F\restriction_{F^{-1}[T]}\colon F^{-1}[T]\bij T}[thm]}\\
\addlinespace[.8ex]
\textbf{Induzierte Potenzmengenabbildung}
\UebFormel{\PowMap{F}(X)\coloneqq F[X].}
Hier ist \(F\colon A\to B\) zunächst nur eine Funktion.
\UebQuellen{\FormulaRefAuto{PowerMapDef}[def]}
&
\textit{Bijektivität wird übertragen und erkannt}
\UebFormel{\begin{aligned}
&F\colon A\bij B\vdash
\PowMap{F}\colon\powerset(A)\bij\powerset(B),\\
&\PowMap{F}\colon\powerset(A)\bij\powerset(B)\\
&\hspace{12mm}\vdash F\colon A\bij B.
\end{aligned}}
\UebQuellen{\FormulaRefAuto{F\colon A\bij B \vdash \PowMap{F}\colon \powerset(A)\bij \powerset(B)}[thm];
\FormulaRefAuto{PowerMapBijectiveReflectsBijectivity}[thm]}\\
\addlinespace[.8ex]
\textbf{Verkleben disjunkter Zweige}
Für \(A_0\cap A_1=\varnothing\), \(B_0\cap B_1=\varnothing\)
und \(f_i\colon A_i\bij B_i\) wird die Fallabbildung verwendet.
\UebQuellen{\FormulaRefAuto{CaseFunctionDef}[def]}
&
\textit{Bijektive Fallabbildung}
\UebFormel{\begin{aligned}
&\CaseFun{A_0}{A_1}{f_0}{f_1}\\
&\hspace{10mm}\colon A_0\cup A_1\bij B_0\cup B_1.
\end{aligned}}
\UebQuellen{\FormulaRefAuto{CaseFunBijectiveDisjointTargets}[thm]}\\
\addlinespace[.8ex]
\textbf{Kernfamilien und Kerntransport}
\UebFormel{\begin{aligned}
&\CoreFamily{P}{A}\coloneqq\\
&\qquad\{H\in\powerset(P\cup A)\mid P\subseteq H\},\\
&\CoreTransport{P}{Q}{f}(H)\coloneqq Q\cup f[H\setminus P].
\end{aligned}}
Für den Transport gilt \(P\cap A=Q\cap B=\varnothing\),
\(f\colon A\to B\), \(H\in\CoreFamily{P}{A}\).
\UebQuellen{\FormulaRefAuto{CoreFamilyDef}[def];
\FormulaRefAuto{CoreTransportDef}[def]}
&
\textit{Bijektiver Transport}
Unter den links genannten Disjunktheitsbedingungen:
\UebFormel{\begin{aligned}
&f\colon A\bij B\vdash\\
&\CoreTransport{P}{Q}{f}\colon
\CoreFamily{P}{A}\bij\CoreFamily{Q}{B}.
\end{aligned}}
\UebQuellen{\FormulaRefAuto{CoreFamilyTransportBijective}[thm]}\\
\addlinespace[.8ex]
\textbf{Schichtfortsetzung}
Für \(g\colon\powerset(D)\to\powerset(E)\),
\(X\in\powerset(A\cup D)\):
\UebFormel{\begin{aligned}
\LayerPowerMap{A}{D}{E}{g}(X)
\coloneqq(X\cap A)\cup g(X\cap D).
\end{aligned}}
\UebQuellen{\FormulaRefAuto{LayerPowerMapDef}[def]}
&
\textit{Bijektive Fortsetzung bei getrennten Schichten}
Für \(A\cap D=A\cap E=\varnothing\) und
\(g\colon\powerset(D)\bij\powerset(E)\):
\UebFormel{\begin{aligned}
\LayerPowerMap{A}{D}{E}{g}\colon
\powerset(A\cup D)\bij\powerset(A\cup E).
\end{aligned}}
Falls außerdem \(g(\varnothing)=\varnothing\), ist auch die
Einschränkung auf die nichtleeren Teilmengen bijektiv.
\UebQuellen{\FormulaRefAuto{LayerPowerMapBijective}[thm];
\FormulaRefAuto{LayerPowerMapNonemptyBijective}[thm]}\\
\addlinespace[.8ex]
\textbf{Cantor--Schröder--Bernstein}
Für \(F\colon A\inj B,\ G\colon B\inj A\) ist
\(\CBPart F G\) die Menge aller \(x\in A\), die zu jedem
\(X\subseteq A\) mit
\UebFormel{\begin{aligned}
&A\setminus G[B]\subseteq X,\\
&G[F[X]]\subseteq X
\end{aligned}}
gehören.
\UebQuellen{\FormulaRefAuto{CantorBernsteinPartDef}[def]}
&
\textit{Fixpunkt und Zusammenfügung}
Mit \(C=\CBPart F G\) gilt
\UebFormel{C=(A\setminus G[B])\cup G[F[C]].}
Auf \(C\) wird \(F\) verwendet, auf \(A\setminus C\) der
eindeutige inverse Wert von \(G\). Dies liefert
\UebFormel{\exists H\colon A\bij B.}
\UebQuellen{\FormulaRefAuto{CantorBernsteinPartFixedPoint}[thm];
\FormulaRefAuto{CantorBernsteinFixedInjections}[thm]}
\\ \addlinespace[.8ex]

\textbf{Erweiterung um ein neues Paar}
Für \(F\colon A\bij B\), \(a\notin A\), \(b\notin B\)
wird die Erweiterungsabbildung auf den vergrößerten Zielbereich angewandt.
\UebQuellen{\FormulaRefAuto{ExtensionMapDef}[def]}
&
\textit{Frische Erweiterung einer Bijektion}
\UebFormel{\begin{aligned}
&\ExtMap{\iota_{B,B\cup\{b\}}\circ F}{a}{b}\\
&\hspace{12mm}\colon A\cup\{a\}\bij B\cup\{b\}.
\end{aligned}}
\UebQuellen{\FormulaRefAuto{FreshExtMapBijective}[thm]}\\
\addlinespace[.8ex]
\textbf{Vorgeschriebener Wert}
Eine vorhandene Bijektion kann durch eine Vertauschung im Ziel
an ein ausgewähltes Paar angepasst werden.
&
\textit{Punktierte Bijektion}
\UebFormel{\begin{aligned}
&\exists u\,(u\colon A\bij B)\dsep a\in A\dsep b\in B\\
&\hspace{6mm}\vdash\exists f\,(f\colon A\bij B\land f(a)=b).
\end{aligned}}
\UebQuellen{\FormulaRefAuto{PointedBijectionExists}[thm]}\\
\end{BandUebersicht}
\noindent\emph{Lesepfad zu Cantor--Schröder--Bernstein.}
Die \CSBReadingLink{} erläutert den kleinsten abgeschlossenen Bereich,
seine Fixpunktgleichung und die Zusammenfügung der beiden Bijektionen.
Die zugeordneten \CSBProofsLink{} enthalten die ausführlichen
Ableitungen zu denselben Resultaten.
